\section{Explanation-Aware Quantization}
\label{sec:eaq}

The theory of Section~\ref{sec:theory} shows that the attribution distortion is
governed by \emph{layer-wise} quantities: the error injected into layer $l$ is
amplified by that layer's contribution to the amplification factor $\Lambda$.
Uniform quantization ignores this heterogeneity and spends the same precision on
every layer. We now derive the precision allocation that minimises explanation
distortion at a fixed average bit budget.

\subsection{Distortion Model}
Under high-resolution quantization, the error power injected into layer $l$ at
bit-width $b_l$ is $D_l\!\approx\! c\,\sigma_l^2\,2^{-2b_l}$, where
$\sigma_l^2=\tfrac1{n_l}\norm{\bW_l}_F^2$ is the mean weight energy and $c$ a
constant (Lemma~\ref{lem:quant} remark). Propagating this error to the
attribution through the layer sensitivity $\Lambda_l=\prod_{k\ne l}\norm{\bW_k}_2$
(Theorem~\ref{thm:lipschitz}), the expected explanation distortion is, to first
order, additive across layers:
\begin{equation}
\mathcal D(\bm b)\;=\;\sum_{l=1}^{L}\omega_l\,2^{-2b_l},\qquad
\omega_l\;\propto\;\Lambda_l^2\,\sigma_l^2 .
\label{eq:distortion}
\end{equation}
The weight $\omega_l$ is the \emph{explanation sensitivity} of layer $l$: it is
large when the layer both carries energy and is strongly amplified. In practice
we estimate $\omega_l$ by a first-order Taylor relevance
$\hat\omega_l=\E_{\bx}\bigl[\,\lVert\,|\bW_l|\odot|\nabla_{\bW_l}\mathcal L|\,
\rVert_1\bigr]$, which captures both magnitude and loss-sensitivity and requires
only a handful of mini-batches. The model \eqref{eq:distortion} assumes a
\emph{uniform} error-decay exponent (the $2^{-2b_l}$ high-resolution rate) for
every layer; Section~\ref{sec:eaq-general} relaxes this to layer-specific
exponents and shows that the uniform choice is the special case $\beta_l\equiv2$.

\subsection{Optimal Bit Allocation}

\begin{theorem}[Reverse water-filling allocation]
\label{thm:waterfill}
The continuous relaxation
\begin{equation}
\min_{\bm b}\;\sum_{l}\omega_l 2^{-2b_l}\quad\text{s.t.}\quad
\frac1L\sum_l b_l=\bar B
\label{eq:opt}
\end{equation}
has the unique solution
\begin{equation}
\boxed{\;b_l^{\star}=\bar B+\tfrac12\log_2\!\Bigl(\frac{\omega_l}{\bar\omega_G}
\Bigr)\;},\qquad
\bar\omega_G=\Bigl(\prod_{k}\omega_k\Bigr)^{1/L},
\label{eq:alloc}
\end{equation}
where $\bar\omega_G$ is the geometric mean of the sensitivities.
\end{theorem}

\begin{proof}
Form the Lagrangian $\mathcal L(\bm b,\lambda)=\sum_l\omega_l2^{-2b_l}+
\lambda\bigl(\sum_l b_l-L\bar B\bigr)$. Stationarity gives
$\partial\mathcal L/\partial b_l=-2\ln2\,\omega_l2^{-2b_l}+\lambda=0$, hence
$2^{-2b_l}=\lambda/(2\ln2\,\omega_l)$ and
$b_l=\tfrac12\log_2\!\bigl(2\ln2\,\omega_l/\lambda\bigr)$. Imposing the budget
$\sum_l b_l=L\bar B$ determines $\lambda$: summing the last expression and solving,
$\tfrac12\log_2(2\ln2/\lambda)=\bar B-\tfrac1{2L}\sum_k\log_2\omega_k$.
Back-substituting yields
$b_l^\star=\bar B+\tfrac12\bigl(\log_2\omega_l-\tfrac1L\sum_k\log_2\omega_k\bigr)$,
which is \eqref{eq:alloc}. The objective is strictly convex in $\bm b$
(each $2^{-2b_l}$ is convex and they are summed with positive weights), so the
stationary point is the unique global minimiser.
\end{proof}

Equation~\eqref{eq:alloc} is a \emph{reverse water-filling}: layers whose
sensitivity exceeds the geometric mean receive more than $\bar B$ bits, those
below receive fewer, and a one-octave (factor-$2$) change in sensitivity shifts
the allocation by exactly half a bit. This is the explanation-preserving analogue
of rate--distortion bit allocation in transform coding. In practice we clip
$b_l^\star$ to an integer box $[b_{\min},b_{\max}]$ and re-project onto the budget
by a short bisection, which is the box-constrained water-filling solution.

\begin{remark}[Integer-rounding optimality]
\label{rem:integer}
Theorem~\ref{thm:waterfill} solves the continuous relaxation; deployable
bit-widths are integers. Because the objective $\sum_l\omega_l2^{-2b_l}$ is
separable and convex and the budget constraint is linear, the integer minimiser
under the box $[b_{\min},b_{\max}]$ is obtained by a greedy marginal-gain
assignment: starting from $b_{\min}$ everywhere, repeatedly add one bit to the
layer with the largest current marginal distortion reduction
$\omega_l(2^{-2b_l}-2^{-2(b_l+1)})=\tfrac34\omega_l2^{-2b_l}$ until the budget is
met. Since the marginal gains are strictly decreasing in $b_l$, this greedy rule
is exactly optimal over integer allocations (a discrete water-filling), and the
rounding of Algorithm~\ref{alg:eaq} coincides with it up to the re-projection
step; the per-layer rounding gap is at most half a bit and never violates the
budget.
\end{remark}

\subsection{Generalised Layer-Specific Rate--Distortion Model}
\label{sec:eaq-general}
The high-resolution model \eqref{eq:distortion} assigns every layer the same
error-decay exponent $2^{-2b_l}$. Empirically, however, the per-layer distortion
of a real quantizer decays at a layer-dependent rate: different depths,
activation dynamics and tensor shapes yield empirical exponents $\beta_l$ that
depart from the ideal value of $2$, a phenomenon recent mixed-precision
rate--distortion analyses term \emph{distortion-model mismatch}
\cite{ranjan2025mixqvit,li2024rdkv}. Forcing a single exponent can misallocate
the budget---occasionally below the quality of plain uniform quantization. We
therefore generalise \eqref{eq:distortion} to
\begin{equation}
\mathcal D(\bm b)\;=\;\sum_{l=1}^{L}\omega_l\,2^{-\beta_l b_l},
\label{eq:distortion-gen}
\end{equation}
where $\beta_l>0$ is estimated per layer by a short pre-computation---quantizing
layer $l$ at two or three probe bit-widths and regressing $\log\mathcal D_l$ on
$b_l$. Repeating the Lagrangian argument of Theorem~\ref{thm:waterfill} on
\eqref{eq:distortion-gen} with $\partial\mathcal D/\partial b_l=
-\beta_l\ln2\,\omega_l2^{-\beta_l b_l}+\lambda=0$ gives the generalised optimal
allocation
\begin{equation}
\boxed{\;b_l^{\star}=\frac{1}{\beta_l}\Bigl[\log_2\!\bigl(\omega_l\,\beta_l\ln2
\bigr)-\log_2\lambda\Bigr]\;},
\label{eq:alloc-gen}
\end{equation}
with the multiplier $\lambda$ fixed by the budget $\tfrac1L\sum_l b_l=\bar B$
through a one-dimensional bisection. Setting $\beta_l\equiv2$ recovers
\eqref{eq:alloc} exactly, so \eqref{eq:alloc-gen} is a strict generalisation:
layers with a \emph{faster} empirical decay (larger $\beta_l$) need fewer bits to
reach the same distortion and are down-weighted accordingly, which prevents the
mismatch failure mode and aligns the allocation with modern rate--distortion
schedulers \cite{li2024rdkv}. Algorithm~\ref{alg:eaq} uses the uniform form by
default and the generalised form when probe measurements of $\beta_l$ are
available.

\subsection{Algorithm}
Algorithm~\ref{alg:eaq} summarises EAQ. It requires only a forward/backward pass
over a few calibration batches to estimate $\hat\omega_l$, followed by a
closed-form allocation and a single quantization pass---no retraining. Unlike
ECQ$^{\mathrm x}$ \cite{becking2022ecq} and the transformer- and
imitation-learning schemes of Table~\ref{tab:xaicomp}, which spend their
relevance signal on \emph{preserving accuracy}, EAQ's objective
\eqref{eq:distortion} is the \emph{explanation distortion} itself; accuracy gains
(Section~\ref{sec:results}) arise as a by-product of protecting the same
high-sensitivity layers.

\begin{algorithm}[!t]
\caption{Explanation-Aware Quantization (EAQ)}
\label{alg:eaq}
\begin{algorithmic}[1]
\REQUIRE trained parameters $\btheta$, calibration set $\mathcal B$, budget
$\bar B$, box $[b_{\min},b_{\max}]$
\ENSURE mixed-precision parameters $\tilde\btheta$
\STATE $\hat\omega_l\leftarrow\frac1{|\mathcal B|}\sum_{\bx\in\mathcal B}
\lVert\,|\bW_l|\odot|\nabla_{\bW_l}\mathcal L(\bx)|\,\rVert_1$ \quad$\forall l$
\hfill{\small\textit{// estimate layer sensitivities}}
\STATE $b_l\leftarrow\bar B+\tfrac12\bigl(\log_2\hat\omega_l-\tfrac1L\sum_k
\log_2\hat\omega_k\bigr)$ \quad$\forall l$
\hfill{\small\textit{// optimal allocation, Thm.~\ref{thm:waterfill}}}
\STATE $b_l\leftarrow\mathrm{clip}(b_l,b_{\min},b_{\max})$
\hfill{\small\textit{// enforce hardware bit range}}
\WHILE{$|\tfrac1L\sum_l\lfloor b_l\rceil-\bar B|>0$}
  \STATE $b_l\leftarrow\mathrm{clip}\bigl(b_l-(\tfrac1L\sum_k\lfloor b_k\rceil
  -\bar B),\,b_{\min},b_{\max}\bigr)$
  \hfill{\small\textit{// re-project onto budget}}
\ENDWHILE
\STATE $\tilde\bW_l\leftarrow Q_{\lfloor b_l\rceil}(\bW_l)$ \quad$\forall l$
\hfill{\small\textit{// quantize each layer at its bit-width}}
\STATE \textbf{return} $\tilde\btheta$
\end{algorithmic}
\end{algorithm}

\subsection{Computational Complexity}
\label{sec:complexity}
To demonstrate the efficiency of EAQ, we analyse the time and space complexity of its components relative to other quantization schemes in Table~\ref{tab:complexity}.

\begin{table}[!h]
\caption{Computational Complexity Comparison of Quantization Schemes ($P$ parameters, $L$ layers, $N_{\mathrm{cal}}$ calibration samples, $N_{\mathrm{iter}}$ optimization steps)}
\label{tab:complexity}
\centering
\begin{tabular}{lccc}
\hline
\textbf{Method} & \textbf{Time (Sensitivity)} & \textbf{Time (Allocation)} & \textbf{Space Overhead} \\
\hline
Uniform (PTQ) & $O(1)$ & $O(1)$ & $O(1)$ \\
AdaRound \cite{nagel2020adaround} & $O(N_{\mathrm{iter}} \cdot P)$ & $O(1)$ & $O(P)$ \\
BRECQ \cite{li2021brecq} & $O(N_{\mathrm{iter}} \cdot P)$ & $O(1)$ & $O(P)$ \\
ECQ$^{\mathrm{x}}$ \cite{becking2022ecq} & $O(N_{\mathrm{cal}} \cdot P)$ & $O(P \log P)$ & $O(P)$ \\
Mix-QViT \cite{ranjan2025mixqvit} & $O(N_{\mathrm{cal}} \cdot P)$ & $O(L)$ & $O(L)$ \\
\textbf{EAQ (ours)} & $O(N_{\mathrm{cal}} \cdot P)$ & $O(L)$ & $O(L)$ \\
\hline
\end{tabular}
\end{table}

\begin{remark}[Time and space complexity]
\emph{Time.} EAQ adds $O(N_{\mathrm{cal}})$ backward passes for sensitivity estimation, plus $O(L)$ arithmetic for the closed-form allocation \eqref{eq:alloc} (or an $O(L\log\frac1\epsilon)$ bisection for the generalised allocation \eqref{eq:alloc-gen}). This is extremely fast (taking less than a minute on CPU) compared to AdaRound or BRECQ, which run iterative optimization over thousands of steps. During inference, EAQ introduces zero computational overhead (0 additional FLOPs) because it compiles to standard mixed-precision layers.

\emph{Space.} The sensitivity estimator accumulates only one scalar $\hat\omega_l$ per layer, using $O(L)$ additional storage. It never stores second-order Hessians or parameter-level rounding variables, achieving significant memory savings over methods that require $O(P)$ auxiliary parameters.
\end{remark}
